\documentclass{article}
\usepackage{../math_template}

\author{Jonathan Kasongo}
\title{British Math Olympiad 1997 Rd 1 --- P4/5}

\begin{document}
\maketitle

\begin{problem}{British Math Olympiad 1997 Rd 1 --- P4/5}
Let $ABCD$ be a convex quadrilateral. The midpoints of $AB$, $BC$, $CD$ and
$DA$ are $P$, $Q$, $R$, and $S$, respectively. Given that the quadrilateral
$PQRS$ has area 1, prove that the area of the quadrilateral $ABCD$ is 2.
\end{problem}


\begin{solution}{Solution}
\begin{center}
\begin{asy}
/* Geogebra to Asymptote conversion, documentation at artofproblemsolving.com/Wiki go to User:Azjps/geogebra */
import graph; size(140);
real labelscalefactor = 0.5; /* changes label-to-point distance */
pen dps = linewidth(0.7) + fontsize(10); defaultpen(dps); /* default pen style */
pen dotstyle = black; /* point style */
real xmin = -7.587127707435243, xmax = 9.499352361738403, ymin = -6.174428887520138, ymax = 5.739568110129332;  /* image dimensions */
pen xdxdff = rgb(0.49019607843137253,0.49019607843137253,1.); pen ccqqqq = rgb(0.8,0.,0.);
pair A = (-3.,1.), B = (5.,4.), C = (6.,-2.), D = (-1.,-2.), P = (1.,2.5), Q = (5.5,1.), R = (2.5,-2.);
/* draw figures */
draw(A--B,  xdxdff);
draw(B--C,  xdxdff);
draw(C--D,  xdxdff);
draw(D--A,  xdxdff);
draw(P--Q,  linetype("4 4") + ccqqqq);
draw(Q--R,  linetype("4 4") + ccqqqq);
draw(R--(-2.,-0.5),  linetype("4 4") + ccqqqq);
draw((-2.,-0.5)--P,  linetype("4 4") + ccqqqq);
/* dots and labels */
dot(A,dotstyle);
label("$A$", (-2.9465195018775767,1.1388505195191572), NE * labelscalefactor);
dot(B,dotstyle);
label("$B$", (5.058197230918168,4.130646640580687), NE * labelscalefactor);
dot(C,dotstyle);
label("$C$", (6.055462604605346,-1.866242473191535), NE * labelscalefactor);
dot(D,dotstyle);
label("$D$", (-0.9519887545032217,-1.866242473191535), NE * labelscalefactor);
dot(P, dotstyle);
label("$P$", (1.0558388645202956,2.601506400927016), NE * labelscalefactor);
dot(Q,dotstyle);
label("$Q$", (5.550181481937176,1.1122567762208324), NE * labelscalefactor);
dot(R, dotstyle);
label("$R$", (2.558385360875643,-1.8928362164898596), NE * labelscalefactor);
dot((-2.,-0.5),dotstyle);
label("$S$", (-1.9492541281903992,-0.39028972013451363), NE * labelscalefactor);
clip((xmin,ymin)--(xmin,ymax)--(xmax,ymax)--(xmax,ymin)--cycle);
/* end of picture */
\end{asy}
\end{center}

The idea is to use the fact that $\triangle BPQ \sim \triangle ABC$.
This because $\overrightarrow{AB} = 2\overrightarrow{PB}$,
$\overrightarrow{CB} = 2\overrightarrow{QB}$ and thus
$\overrightarrow{AC} = \overrightarrow{AB} + \overrightarrow{BC} =
2(\overrightarrow{PB} + \overrightarrow{BQ}) = 2\overrightarrow{PQ}$. Thus we can conclude
$[ABC] : [PBQ] = 4 : 1$. By symmetrical reasoning we have
$$
\frac{[PBQ]}{[ABC]} = \frac{[SDR]}{[ADC]} = \frac{[QCR]}{[BCD]} =
\frac{[PAS]}{[BAD]} = \frac{1}{4}
$$
and so
\[
\begin{aligned}
[ABCD] &= [ABC] + [ADC] = 4([PBQ] + [ADC])\\
&= [DAB] + [DCB] = 4([SAP] + [RCQ])
\end{aligned}
\]
which gives
\[
\begin{aligned}
2[ABCD] &= 4([PBQ] + [ADC] + [SAP] + [RCQ])\\
[ABCD] &= 2([PBQ] + [ADC] + [SAP] + [RCQ])\\
&=  ([PBQ] + [ADC] + [SAP] + [RCQ]) + [PQRS]\\
\end{aligned}
\]
and so we can conclude that $([PBQ] + [ADC] + [SAP] + [RCQ]) = [PQRS] = 1$.
Now substituting gives that $[ABCD] = 2$. $\blacksquare$
\end{solution}


\end{document}
